\section{Evaluation Prompts}
\label{sec:eval_prompts}


To ensure reproducibility and facilitate future research, we provide here the complete set of prompts used to evaluate our model across all benchmarks. These prompts were consistently applied during inference to maintain fairness and comparability.

\subsection{STEM \& Puzzle}
\label{sec:stem}

\begin{tcolorbox}[
  title=MMMU,
  fonttitle=\bfseries,
  colframe=black,
  colback=white,
  toptitle=1mm,
  bottomtitle=1mm,
  top=2mm,
  verbatim,
]
<image>

Question: \{question\}

Options:

\{options\}

Please select the correct answer from the options above.
\end{tcolorbox}

\begin{tcolorbox}[
  title=MMMUPro\_Standard,
  fonttitle=\bfseries,
  colframe=black,
  colback=white,
  toptitle=1mm,
  bottomtitle=1mm,
  top=2mm,
  verbatim,
]
<image>

\{question\}

\{options\}

Please select the correct answer from the options.
\end{tcolorbox}

\begin{tcolorbox}[
  title=MMMUPro\_Vision,
  fonttitle=\bfseries,
  colframe=black,
  colback=white,
  toptitle=1mm,
  bottomtitle=1mm,
  top=2mm,
  verbatim,
]
<image>

Identify the problem and solve it. Think step by step before answering.
\end{tcolorbox}

\begin{tcolorbox}[
  title=MathVista,
  title={MathVista | MathVision | MathVerse | LogicVista},
  fonttitle=\bfseries,
  colframe=black,
  colback=white,
  toptitle=1mm,
  bottomtitle=1mm,
  top=2mm,
  verbatim,
]
<image>

\{question\}
\end{tcolorbox}

\begin{tcolorbox}[
  title=We-Math,
  fonttitle=\bfseries,
  colframe=black,
  colback=white,
  toptitle=1mm,
  bottomtitle=1mm,
  top=2mm,
  verbatim,
]
<image>

Now, we require you to solve a multiple-choice math question. Please briefly describe your thought process and provide the final answer(option).

Question: \{question\}

Option: \{options\}

Regarding the format, please answer following the template below, and be sure to include two <> symbols:

<Thought process>: <<your thought process>> <Answer>: <<your option>>
\end{tcolorbox}

\begin{tcolorbox}[
  title=ZeroBench,
  fonttitle=\bfseries,
  colframe=black,
  colback=white,
  toptitle=1mm,
  bottomtitle=1mm,
  top=2mm,
  verbatim,
]
<image>

\{question\}

Let's think step by step and give the final answer in curly braces, like this: \{final answer\}
\end{tcolorbox}

\begin{tcolorbox}[
  title=DynaMath,
  fonttitle=\bfseries,
  colframe=black,
  colback=white,
  toptitle=1mm,
  bottomtitle=1mm,
  top=2mm,
  verbatim,
]
<image>

\#\# Question

\{question\}

\#\# Answer Instruction: Please provide an answer to the question outlined above. Your response should adhere to the following JSON format, which includes two keys: 'solution' and 'short answer'. The 'solution' key can contain detailed steps needed to solve the question, and the 'short answer' key should provide a concise response. 


Example of expected JSON response format:

\{

    "solution": "[Detailed step-by-step explanation]",
    
    "short answer": "[Concise Answer]"
    
\}
\end{tcolorbox}

\begin{tcolorbox}[
  title=VLMBlind,
  fonttitle=\bfseries,
  colframe=black,
  colback=white,
  toptitle=1mm,
  bottomtitle=1mm,
  top=2mm,
  verbatim,
]
<image>

Question: \{question\}
\end{tcolorbox}

\begin{tcolorbox}[
  title=VisuLogic,
  fonttitle=\bfseries,
  colframe=black,
  colback=white,
  toptitle=1mm,
  bottomtitle=1mm,
  top=2mm,
  verbatim,
]
<image>

\{question\}

Solve the complex visual logical reasoning problem through step-by-step reasoning. Think about the reasoning process first and answer the question following this format: Answer:{/}/boxed\{\$LETTER\}
\end{tcolorbox}

\begin{tcolorbox}[
  title=VisualPuzzles-Direct,
  fonttitle=\bfseries,
  colframe=black,
  colback=white,
  toptitle=1mm,
  bottomtitle=1mm,
  top=2mm,
  verbatim,
]
<image>

Question: \{question\}

Options:

\{options\}

Answer the question with the option's letter from the given choices directly.
\end{tcolorbox}

\begin{tcolorbox}[
  title=VisualPuzzles-CoT,
  fonttitle=\bfseries,
  colframe=black,
  colback=white,
  toptitle=1mm,
  bottomtitle=1mm,
  top=2mm,
  verbatim,
]
<image>

Question: \{question\}

Options:

\{options\}

Solve the multiple-choice question and then answer with the option letter from the given choices. The last line of your response should be of the following format: 'Answer: \$LETTER' (without quotes), where LETTER is one of the options. Think step by step before answering.
\end{tcolorbox}

\subsection{GeneralVQA}
\label{sec:vqa}

\begin{tcolorbox}[
  title={MMBench | RealWorldQA | MMStar},
  fonttitle=\bfseries,
  colframe=black,
  colback=white,
  toptitle=1mm,
  bottomtitle=1mm,
  top=2mm,
  verbatim,
]
<image>

Question: \{question\}

Options:

\{options\}

Please select the correct answer from the options above.
\end{tcolorbox}

\begin{tcolorbox}[
  title=SimpleVQA,
  fonttitle=\bfseries,
  colframe=black,
  colback=white,
  toptitle=1mm,
  bottomtitle=1mm,
  top=2mm,
  verbatim,
]
<image>

\{question\}
\end{tcolorbox}

\subsection{Alignment}
\label{sec:alignment}

\begin{tcolorbox}[
  title={HallusionBench | MM\_MT\_Bench | MIA-Bench},
  fonttitle=\bfseries,
  colframe=black,
  colback=white,
  toptitle=1mm,
  bottomtitle=1mm,
  top=2mm,
  verbatim,
]
<image>

\{question\}
\end{tcolorbox}

\subsection{Document-Understanding}
\label{sec:doc}

\begin{tcolorbox}[
  title=MMLongBench-Doc,
  fonttitle=\bfseries,
  colframe=black,
  colback=white,
  toptitle=1mm,
  bottomtitle=1mm,
  top=2mm,
  verbatim,
]
<image\_1>

<image\_2>

...

<image\_n>

\{question\}
\end{tcolorbox}

\begin{tcolorbox}[
  title={DocVQA | InfoVQA | ChartQA\_TEST},
  fonttitle=\bfseries,
  colframe=black,
  colback=white,
  toptitle=1mm,
  bottomtitle=1mm,
  top=2mm,
  verbatim,
]
<image>

\{question\}

Answer the question using a single word or phrase.
\end{tcolorbox}

\begin{tcolorbox}[
  title=AI2D,
  fonttitle=\bfseries,
  colframe=black,
  colback=white,
  toptitle=1mm,
  bottomtitle=1mm,
  top=2mm,
  verbatim,
]
<image>

Question: \{question\}

Options:

\{options\}

Please select the correct answer from the options above.
\end{tcolorbox}

\begin{tcolorbox}[
  title={OCRBench | OCRBench\_v2 | CC-OCR | CharXiv},
  fonttitle=\bfseries,
  colframe=black,
  colback=white,
  toptitle=1mm,
  bottomtitle=1mm,
  top=2mm,
  verbatim,
]
<image>

\{question\}
\end{tcolorbox}

\begin{tcolorbox}[
  title=OmniDocBench,
  fonttitle=\bfseries,
  colframe=black,
  colback=white,
  toptitle=1mm,
  bottomtitle=1mm,
  top=2mm,
  verbatim,
]
<image>

You are an AI assistant specialized in converting PDF images to Markdown format. Please follow these instructions for the conversion:

1. Text Processing:
- Accurately recognize all text content in the PDF image without guessing or inferring.
- Convert the recognized text into Markdown format.
- Maintain the original document structure, including headings, paragraphs, lists, etc.

\begin{verbatim}
2. Mathematical Formula Processing:
- Convert all mathematical formulas to LaTeX format.
- Enclose inline formulas with \( \). For example: This is an inline formula \( E = mc^2 \)
- Enclose block formulas with \[ \]. For example: \[ \frac{-b \pm \sqrt{b^2 - 4ac}}{2a} \]
\end{verbatim}

3. Table Processing:
- Convert tables to HTML format.
- Wrap the entire table with <table> and </table>.

4. Figure Handling:
- Ignore figures in the PDF image. Do not attempt to describe or convert images.

5. Output Format:
- Ensure the output Markdown document has a clear structure with appropriate line breaks between elements.
- For complex layouts, try to maintain the original document's structure and format as closely as possible.

Please strictly follow these guidelines to ensure accuracy and consistency in the conversion. Your task is to accurately convert the content of the PDF image into Markdown format without adding any extra explanations or comments.
\end{tcolorbox}

\subsection{2D/3D Grounding}
\label{sec:grounding}

\begin{tcolorbox}[
  title=RefCOCO,
  fonttitle=\bfseries,
  colframe=black,
  colback=white,
  toptitle=1mm,
  bottomtitle=1mm,
  top=2mm,
  verbatim,
]
<image>

Locate every object that matches the description "\{ref\_sentence\}" in the image. Report bbox coordinates in JSON format.
\end{tcolorbox}

\begin{tcolorbox}[
  title=CountBench,
  fonttitle=\bfseries,
  colframe=black,
  colback=white,
  toptitle=1mm,
  bottomtitle=1mm,
  top=2mm,
  verbatim,
]
<image>

Question: \{question\}

Options:

\{options\}

Please select the correct answer from the options above.
\end{tcolorbox}

\begin{tcolorbox}[
  title=ODinW-13,
  fonttitle=\bfseries,
  colframe=black,
  colback=white,
  toptitle=1mm,
  bottomtitle=1mm,
  top=2mm,
  verbatim,
]
<image>

Locate every instance that belongs to the following categories: \'\{obj\_names\}\'. Report bbox coordinates in JSON format.
\end{tcolorbox}

\begin{tcolorbox}[
  title={ARKitScenes | Hypersim  | SUNRGBD},
  fonttitle=\bfseries,
  colframe=black,
  colback=white,
  toptitle=1mm,
  bottomtitle=1mm,
  top=2mm,
  verbatim,
]
<image>

Locate the \{class\_name \} in the provided image and output their positions and dimensions using 3D bounding boxes. The results must be in the JSON format: \texttt{[{"bbox\_3d":[x\_center, y\_center, z\_center, x\_size, y\_size, z\_size, roll, pitch, yaw],"label":"category"}]}. 
\end{tcolorbox}

\subsection{Embodied/Spatial
Understanding}
\label{sec:embodied}

\begin{tcolorbox}[
  title=ERQA,
  fonttitle=\bfseries,
  colframe=black,
  colback=white,
  toptitle=1mm,
  bottomtitle=1mm,
  top=2mm,
  verbatim,
]
<image\_1>

<image\_2>

...

<image\_n>

\{question\}
\end{tcolorbox}


\begin{tcolorbox}[
  enhanced,
  title=VSI-Bench,
  fonttitle=\bfseries,
  colframe=black,
  colback=white,
  toptitle=1mm,
  bottomtitle=1mm,
  top=1mm,
  bottom=1mm,
  left=1mm,
  right=1mm,
  middle=1mm,
  segmentation style={solid, draw=black, line width=0.4pt}
]
\textbf{multiple-choice:}
\begin{Verbatim}
<video>
These are frames of a video.
{question}
Options:
{options}
Answer with the option's letter from the given choices directly.
\end{Verbatim}

\tcblower

\textbf{open-ended:}
\begin{Verbatim}
<video>
These are frames of a video.
{question}
Please answer the question using a single word or phrase.
\end{Verbatim}
\end{tcolorbox}

\begin{tcolorbox}[
  title=EmbSpatialBench,
  fonttitle=\bfseries,
  colframe=black,
  colback=white,
  toptitle=1mm,
  bottomtitle=1mm,
  top=2mm,
  verbatim,
]
<image>

\{question\}
\end{tcolorbox}

\begin{tcolorbox}[
  title=RoboSpatialHome,
  fonttitle=\bfseries,
  colframe=black,
  colback=white,
  toptitle=1mm,
  bottomtitle=1mm,
  top=2mm,
  verbatim,
]
<image>

Locate \{object\_name\} in this image. Output the point coordinates in JSON format.

For example:


[

    \{"point\_2d": [x, y], "label": "point\_1"\}
    
]
\end{tcolorbox}

\begin{tcolorbox}[
  title=RefSpatialBench,
  fonttitle=\bfseries,
  colframe=black,
  colback=white,
  toptitle=1mm,
  bottomtitle=1mm,
  top=2mm,
  verbatim,
]
<image>

\{question\} Output the point coordinates in JSON format.

For example:


[

     \{"point\_2d": [x, y], "label": "point\_1"\}
     
]
\end{tcolorbox}

\subsection{Multi-Image}
\label{sec:multi-img}

\begin{tcolorbox}[
  title=BLINK,
  fonttitle=\bfseries,
  colframe=black,
  colback=white,
  toptitle=1mm,
  bottomtitle=1mm,
  top=2mm,
  verbatim,
]
<image>

Question: \{question\}

Options:

\{options\}

Please select the correct answer from the options above.
\end{tcolorbox}

\begin{tcolorbox}[
  title=MUIRBENCH,
  fonttitle=\bfseries,
  colframe=black,
  colback=white,
  toptitle=1mm,
  bottomtitle=1mm,
  top=2mm,
  verbatim,
]
<image\_1>

<text\_1>

<image\_2>

<text\_2>

...

<image\_n>

<text\_n>

Answer with the option's letter from the given choices directly.
\end{tcolorbox}

\subsection{Video Understanding}
\label{sec:video}

\begin{tcolorbox}[
  title={MVBench | VideoMME | MLVU | LVBench - For instruct models},
  fonttitle=\bfseries,
  colframe=black,
  colback=white,
  toptitle=1mm,
  bottomtitle=1mm,
  top=2mm,
  verbatim,
]
<video>

Select the best answer to the following multiple-choice question based on the video. 

Respond with only the letter (A, B, C, or D) of the correct option.

Question: \{question\} Possible answer choices:

\{options\}

The best answer is:
\end{tcolorbox}

\begin{tcolorbox}[
  title={MVBench | VideoMME | MLVU | LVBench - For thinking models},
  fonttitle=\bfseries,
  colframe=black,
  colback=white,
  toptitle=1mm,
  bottomtitle=1mm,
  top=2mm,
  verbatim,
]
<video>

Select the best answer to the following multiple-choice question based on the video. Respond with only the letter (A, B, C, or D) of the correct option.

Question: \{question\}

\{options\}

Please reason step-by-step, identify relevant visual content, analyze key timestamps and clues, and then provide the final answer.
\end{tcolorbox}

\begin{tcolorbox}[
  title=Charades-STA,
  fonttitle=\bfseries,
  colframe=black,
  colback=white,
  toptitle=1mm,
  bottomtitle=1mm,
  top=2mm,
  verbatim,
]
<video>

Give you a textual query: \{query\_text\}

When does the described content occur in the video? 

Please return the timestamp in seconds.
\end{tcolorbox}

\begin{tcolorbox}[
  enhanced,
  title=VideoMMMU,
  fonttitle=\bfseries,
  colframe=black,
  colback=white,
  toptitle=1mm,
  bottomtitle=1mm,
  top=1mm,
  bottom=1mm,
  left=1mm,
  right=1mm,
  middle=1mm,
  segmentation style={solid, draw=black, line width=0.4pt}
]
% --- 第一部分 ---
\textbf{Perception \& Comprehension:}
\begin{Verbatim}
<video>
{question}
{options}
Please ignore the Quiz question in last frame of the video.
\end{Verbatim}

\tcblower % <--- 第一个分割用 \tcblower，这会画出第一条线

% --- 第二部分 ---
\textbf{Adaptation-multiple-choice:}
\begin{Verbatim}
<video>
<image>
You should watch and learn the video content. Then apply what you learned to answer the 
following multi-choice question. The image for this question is at the end of the video.
{question}
{options}
\end{Verbatim}

\tcbline % <--- 用 \tcbline 手动绘制第二条线

% --- 第三部分 ---
\textbf{Adaptation-open-ended:}
\begin{Verbatim}
<video>
<image>
You should watch and learn the video content. Then apply what you learned to answer the 
following open-ended question. The image for this question is at the end of the video.
{question}
\end{Verbatim}
\end{tcolorbox}

\begin{tcolorbox}[
  enhanced,
  title=MMVU,
  fonttitle=\bfseries,
  colframe=black,
  colback=white,
  toptitle=1mm,
  bottomtitle=1mm,
  top=1mm,
  bottom=1mm,
  left=1mm,
  right=1mm,
  middle=1mm,
  segmentation style={solid, draw=black, line width=0.4pt},
]

\textbf{multiple-choice:}

\begin{Verbatim}[breaklines=true, breakanywhere=true, breaksymbolleft={}]
<video>
{question}
{options}
Visual Information: processed video
Answer the given multiple-choice question step by step. Begin by explaining your reasoning process clearly. Conclude by stating the final answer using the following format: "Therefore, the final answer is: $LETTER" (without quotes), where $LETTER is one of the options. Think step by step before answering.
\end{Verbatim}
\tcbline

\textbf{open-ended:}
\begin{Verbatim}[breaklines=true, breakanywhere=true, breaksymbolleft={}]
<video>
{question}
Visual Information: processed video
Answer the given question step by step. Begin by explaining your reasoning process clearly. 
Conclude by stating the final answer using the following format: "Therefore, the final answer is: "Answer: $ANSWER" (without quotes), where $ANSWER is the final answer of the question. Think step by step before answering.
\end{Verbatim}

\end{tcolorbox}

\subsection{Perception with Tool}
\label{sec:perception_w_tool}

\begin{tcolorbox}[
  title=V*,
  fonttitle=\bfseries,
  colframe=black,
  colback=white,
  toptitle=1mm,
  bottomtitle=1mm,
  top=2mm,
  verbatim,
]

Your role is that of a research assistant specializing in visual information. Answer questions about images by looking at them closely and then using research tools. Please follow this structured thinking process and show your work.

\

Start an iterative loop for each question:

\

- **First, look closely:** Begin with a detailed description of the image, paying attention to the user's question. List what you can tell just by looking, and what you'll need to look up.

- **Next, find information:** Use a tool to research the things you need to find out.

- **Then, review the findings:** Carefully analyze what the tool tells you and decide on your next action.

\

Continue this loop until your research is complete.

\

To finish, bring everything together in a clear, synthesized answer that fully responds to the user's question.

\

\#Tools

\ 

You may call one or more functions to assist with the user query.

\ 

You are provided with function signatures within \texttt{\textless tools\textgreater}\texttt{\textless /tools\textgreater} XML tags:
\par\noindent\texttt{\textless tools\textgreater}

\texttt{\{}
\texttt{"type":"function",}
\texttt{"function":\ \{"name":\ "image\_zoom\_in\_tool",\ "description":\ "Zoom in on a specific region of an image by cropping it based on a bounding box (bbox) and an optional object label",\ "arguments":\ \{"type":\ "object",\ "properties":\ \{"bbox\_2d":\ \{"type":\ "array",\ "items":\ \{"type":\ "number"\},\ "minItems":\ 4,\ "maxItems":\ 4,\ "description":\ "The bounding\ box\ of\ the\ region to zoom in, as [x1, y1, x2, y2], where (x1, y1) is the top-left corner and (x2, y2) is the bottom-right corner"\}, "label": \{"type": "string", "description": "The name or label of the object in the specified bounding box"\}, "img\_idx": \{"type": "number", "description": "The index of the zoomed-in image (starting from 0)"\}\}, "required": ["bbox\_2d", "label", "img\_idx"]}\}\}\}

\texttt{\textless /tools\textgreater}

For each function call, return a JSON object with function name and arguments within \texttt{\textless tool\_call\textgreater}\texttt{\textless /tool\_call\textgreater} XML tags:
\par\noindent\texttt{\textless tool\_call\textgreater}\\
\texttt{\{\{"name": \textless function-name\textgreater , "arguments": \textless args-json-object\textgreater \}\}}\\
\texttt{\textless /tool\_call\textgreater}

<image>

\{question\}

\end{tcolorbox}

\begin{tcolorbox}[
  title={HRBench4K | HRBench8K},
  fonttitle=\bfseries,
  colframe=black,
  colback=white,
  toptitle=1mm,
  bottomtitle=1mm,
  top=2mm,
  verbatim,
]

Your role is that of a research assistant specializing in visual information. Answer questions about images by looking at them closely and then using research tools. Please follow this structured thinking process and show your work.

\

Start an iterative loop for each question:

\

- **First, look closely:** Begin with a detailed description of the image, paying attention to the user's question. List what you can tell just by looking, and what you'll need to look up.

- **Next, find information:** Use a tool to research the things you need to find out.

- **Then, review the findings:** Carefully analyze what the tool tells you and decide on your next action.

\

Continue this loop until your research is complete.

\

To finish, bring everything together in a clear, synthesized answer that fully responds to the user's question.

\

\#Tools

\ 

You may call one or more functions to assist with the user query.

\ 

You are provided with function signatures within \texttt{\textless tools\textgreater}\texttt{\textless /tools\textgreater} XML tags:
\par\noindent\texttt{\textless tools\textgreater}

\texttt{\{}
\texttt{"type":"function",}
\texttt{"function":\ \{"name":\ "image\_zoom\_in\_tool",\ "description":\ "Zoom in on a specific region of an image by cropping it based on a bounding box (bbox) and an optional object label",\ "arguments":\ \{"type":\ "object",\ "properties":\ \{"bbox\_2d":\ \{"type":\ "array",\ "items":\ \{"type":\ "number"\},\ "minItems":\ 4,\ "maxItems":\ 4,\ "description":\ "The bounding\ box\ of\ the\ region to zoom in, as [x1, y1, x2, y2], where (x1, y1) is the top-left corner and (x2, y2) is the bottom-right corner"\}, "label": \{"type": "string", "description": "The name or label of the object in the specified bounding box"\}, "img\_idx": \{"type": "number", "description": "The index of the zoomed-in image (starting from 0)"\}\}, "required": ["bbox\_2d", "label", "img\_idx"]}\}\}\}

\texttt{\textless /tools\textgreater}

For each function call, return a JSON object with function name and arguments within \texttt{\textless tool\_call\textgreater}\texttt{\textless /tool\_call\textgreater} XML tags:
\par\noindent\texttt{\textless tool\_call\textgreater}\\
\texttt{\{\{"name": \textless function-name\textgreater , "arguments": \textless args-json-object\textgreater \}\}}\\
\texttt{\textless /tool\_call\textgreater}

<image>

\{question\}

\{options\}

\end{tcolorbox}


\subsection{Coding}


\begin{tcolorbox}[
  title=Design2Code (Generation),
  fonttitle=\bfseries,
  colframe=black,
  colback=white,
  toptitle=1mm,
  bottomtitle=1mm,
  top=2mm,
  verbatim,
]
<image>

You are an expert web developer who specializes in HTML and CSS.
A user will provide you with a screenshot of a webpage.
You need to return a single HTML file that uses HTML and CSS to reproduce the given website.
Include all CSS code in the HTML file itself.
If it involves any images, use "rick.jpg" as the placeholder.
Some images on the webpage are replaced with a blue rectangle as the placeholder, and use "rick.jpg" for those as well.
Do not hallucinate any dependencies on external files. You do not need to include JavaScript scripts for dynamic interactions.
Pay attention to things like size, text, position, and color of all the elements, as well as the overall layout.
Respond with the content of the HTML+CSS file:
\end{tcolorbox}

\begin{tcolorbox}[
  title=Design2Code (GPT-o4-mini Evaluation),
  fonttitle=\bfseries,
  colframe=black,
  colback=white,
  toptitle=1mm,
  bottomtitle=1mm,
  top=2mm,
  verbatim,
]
I will give you two images. The first is the reference, and the second is generated from the first via code rendering. Please rate their similarity from 0-100, where 0 means completely different and 100 means identical. Provide the score inside a LaTeX \boxed{} and briefly explain your reasoning.

<reference\_image>

<generated\_image>
\end{tcolorbox}

\subsection{Agent}
\begin{tcolorbox}[
  title={Screenspot | Screenspot-Pro | OSWorld-G},
  fonttitle=\bfseries,
  colframe=black,
  colback=white,
  toptitle=1mm,
  bottomtitle=1mm,
  top=2mm,
  breakable
]
\textbf{Tools}

You may call one or more functions to assist with the user query.

You are provided with function signatures within \texttt{\textless tools\textgreater} \dots{} \texttt{\textless /tools\textgreater} XML tags:
\par\noindent\texttt{\textless tools\textgreater}
\texttt{\{}
\texttt{"name":"computer\_use",}
\texttt{"description": "Use a mouse to interact with a computer. The screen's resolution is \textless display\_width\_px\textgreater x \textless display\_height\_px\textgreater ."}
\texttt{"notes": "Click with the cursor tip centered on targets; avoid edges unless asked. Do not use other tools (type, key, scroll, left\_click\_drag). Only left\_click and mouse\_move are allowed. If you can't find the element, terminate and report failure.",}
\texttt{"parameters":\{}
\texttt{\ \ "type":"object",}
\texttt{\ \ "required":["action"],}
\texttt{\ \ "properties":\{}
\texttt{\ \ \ \ "action":\{}
\texttt{\ \ \ \ \ \ "type":"string",}
\texttt{\ \ \ \ \ \ "enum":["mouse\_move","left\_click"],}
\texttt{\ \ \ \ \ \ "description":"The action to perform."}
\texttt{\ \ \ \ \},}
\texttt{\ \ \ \ "coordinate":\{}
\texttt{\ \ \ \ \ \ "type":"array",}
\texttt{\ \ \ \ \ \ "description":"(x, y): pixels from left/top. Required for action=mouse\_move and action=left\_click."}
\texttt{\}}
\texttt{\}}
\texttt{\}}
\texttt{\}}

\texttt{\textless /tools\textgreater}

For each function call, return a JSON object with function name and arguments within \texttt{\textless tool\_call\textgreater} \dots{} \texttt{\textless /tool\_call\textgreater} XML tags:
\par\noindent\texttt{\textless tool\_call\textgreater}\\
\texttt{\{\{"name": \textless function-name\textgreater , "arguments": \textless args-json-object\textgreater \}\}}\\
\texttt{\textless /tool\_call\textgreater}

Additionally, if you think the task is infeasible (e.g., the task is not related to the image), return:
\par\noindent\texttt{\textless tool\_call\textgreater}\\
\texttt{\{"name": "computer\_use", "arguments": \{"action": "terminate", "status": "failure"\}\}}\\
\texttt{\textless /tool\_call\textgreater}
\end{tcolorbox}

